\section{Evaluation}
\label{sec:validation}

Our evaluation assigns different claims to four evidence layers.
Lean checks the mathematical results; an executable cross-check tests the
compiler; litmus manifests exercise distinctions that simpler policies miss;
and public plus private traces characterize workload shape and adapter
observability.
No observational trace is labeled safe or unsafe, and no trace statistic is
used to establish either inclusion involving $\mathsf{Actual}$.

\subsection{Mechanized theorem spine}

The Lean~4 development covers four linked layers.
The typed-history module derives may and required families for all six
operators, proves $\mathsf{Required}\subseteq\mathsf{Candidate}$, and checks
the four semantic outcomes.
The durable-prefix module proves that the admitted residual is prefix-safe and
greatest within the fixed supplied candidate.
The controller-cover module formalizes overlapping cell/controller access,
local families, co-liveness, and the raw physical product.
It proves the readiness and obstruction equivalences, transports an admitted
prefix certificate to an actual subfamily under the explicit refinement
premises, and constructs the older functional partition as a canonical special
case.
The operational commitment module supplies append-only receipt, phase
monotonicity, retry binding, and exact fresh-event accounting.

The audit scans all project Lean sources for placeholders and project axioms,
checks a frozen paper-theorem list, builds the library, prints every listed
kernel dependency, rejects dependencies outside
\texttt{propext}, \texttt{Quot.sound}, and
\texttt{Classical.choice}, and replays the root with a fresh kernel checker.
The U(2,3) fixture is part of this replay: every pair of three cells is
admitted, the triple is a minimal nonface, and two controllers with admitted
proper local pieces physically realize that forbidden triple.

This proves a finite manifest-relative theorem.
It does not prove that a Claude or Codex adapter reports complete contracts,
that its declared controller product contains every actual behavior, or that
an external effect executes exactly once.

\subsection{Compiler--verifier cross-check}

The dependency-free Python artifact contains two semantic implementations.
The compiler uses symbolic cell sets; the verifier does not import it and
reconstructs the request with integer bit masks.  The artifact suite has 55
tests: 31 focus on history admission, and 24 exercise the supporting authority
continuity layer.  The focused tests cover all operators and outcomes,
readiness grades, both identity normalizations, every failure class, malformed
evidence, resource caps, determinism, and tampering.  The end-to-end positive
seal still sets $\mathsf{effect\_authorizes}=\mathsf{false}$.

An exhaustive oracle enumerates all 19 downward-closed families over three
cells.  Across every full-support candidate, every required subfamily, and
every physical family, it checks 2,166 refinement instances.
The compiler's five deployment grades agree with direct set differences in
every instance, and the independent verifier accepts every untampered result.
This enumeration checks the finite algorithm and its independent
reimplementation; it is not a proof that the JSON parser refines Lean.

A separate bounded characterization exercises the production CLIs on
threshold families through 12 cells.  Their configuration/nonface counts are
$11/4$, $42/15$, $163/56$, $638/210$, and $2{,}510/792$ for
$4,6,8,10,12$ cells.  With six cells and one through four overlapping co-live
controllers, the final physical family stays at 64 configurations while the
analytically reconstructed pair-expansion iterations rise from 64 to 70,592.
Both implementations reject 13 source cells, and a five-controller instance
whose uncapped expansion would make 202,688 iterations stops at iteration
200,001.  These synthetic counts demonstrate bounded completion and the two
exercised fail-closed limits; they are neither latency nor deployability
evidence, and the other declared limits remain untested.

\subsection{Distinguishing litmus manifests}

Four litmus manifests separate the failure classes.  A prefix-spent restored
cell exposed by a stale controller yields \textsc{OutsideSupport}; one locally
unsafe controller yields \textsc{LocalOverpermission}; three locally safe
one-cell controllers that compose the U(2,3) triple yield a higher-order
\textsc{GateCut}; and cloned controller anchors that product-compose one
exclusive gate yield \textsc{GateClone}.  Aliasing handles to one cell does
not repair controller cloning.  Conversely, multiple controllers are harmless
for a full powerset future, so the rule does not require one global monitor.

\subsection{Public workload evidence}

TraceLab collects real Claude Code and Codex workload telemetry and reports
long autonomous loops and heterogeneous tool behavior
~\cite{zhu2026tracelab}.
We verified its v0.0.2 release of 8,058 sessions and 743,819 normalized tool
calls against the published artifact digest~\cite{tracelab2026release}.
Its normalized schema provides ordinary call/result correlation but not the
joint typed history edge, durable receipt prefix, cell-alias evidence,
controller identity and co-liveness, or a correlated leaf future contract.
The missing fields make history-admission prevalence unidentifiable; they do
not imply that failures are rare.

\subsection{The paper-formation lineage}

We also analyze the paper-formation trajectory as one author-operated
longitudinal case.  A fixed-cutoff recursive lineage contains over 80 rollout
files, roughly 39,000 native rows, and roughly 6,000 tool calls after excluding
over 200,000 copied-prefix rows.  It exhibits branching, inherited history,
compaction, aborted turns, lifecycle events, workspace mutation, and
call/result correlation.  Yet it does not jointly expose complete leaf
contracts, cell aliases, controller origin/co-liveness, effect binding, and
durable receipt phase.  It therefore motivates adapter fields and parser
tests, but cannot estimate unsafe-restore prevalence or validate admission.
The source-grounded extraction and its limitations appear in
\cref{app:private-trace}.

\subsection{Privacy and ethics}

The case recruited or manipulated no users, but histories may contain private
workspace and incidental third-party text.  Extraction emits only path-free
allowlisted aggregates and private keyed commitments; raw histories, prompts,
commands, paths, stable identifiers, the key, and the anonymous private summary
are not released.  We make no IRB approval or exemption claim; further
minimization details appear in \cref{app:private-trace}.
